\documentclass[conference]{IEEEtran}

\usepackage[bahasa,english]{babel}
\usepackage{graphicx}
\usepackage{cite}
\usepackage{amsmath,amssymb}
\usepackage{hyperref}

\hypersetup{
    colorlinks=true,
    linkcolor=black,
    citecolor=black,
    urlcolor=blue
}

\begin{document}

\title{Judul Paper IEEE}
\author{
    \IEEEauthorblockN{Penulis Pertama\IEEEauthorrefmark{1},
                       Penulis Kedua\IEEEauthorrefmark{2}}
    \IEEEauthorblockA{\IEEEauthorrefmark{1}Prodi, Fakultas, Universitas \\
                      email@example.com}
    \IEEEauthorblockA{\IEEEauthorrefmark{2}Prodi, Fakultas, Universitas \\
                      email2@example.com}
}

\maketitle

\begin{abstract}
Abstrak paper ini.
\end{abstract}

\begin{IEEEkeywords}
kata kunci, dipisah, koma
\end{IEEEkeywords}

\section{Pendahuluan}
\label{sec:pendahuluan}

Tulis pendahuluan di sini.

\section{Metode Penelitian}
\label{sec:metode}

Tulis metode di sini.

\section{Hasil dan Pembahasan}
\label{sec:hasil}

Tulis hasil di sini.

\section{Kesimpulan}
\label{sec:kesimpulan}

Tulis kesimpulan di sini.

\bibliographystyle{IEEEtran}
\bibliography{references}

\end{document}
